% \todo{Results of running some different Rizzo programs (with valgrind) for extended period/number of iterations. Goal is to convince no space is leaks happen when the program is run.}

For this part of the evaluation our focus will be on space leaks. For this purpose, we will be using Valgrind \cite{valgrind}, which is a framework for building dynamic analysis tools.
Valgrind comes with a tool, \textit{memcheck}, for detecting memory errors such as reading uninitialised values, use-after-free, and most importantly memory leak detection. Valgrind cant prove we have no space leaks, but it can give us confidence that we don't have any obvious memory that we are losing track of.

To get information on how many inputs the program has processed, we instrument the runtime to save the number of steps taken and the average time taken per step in microseconds. 
This information is logged at program exit when the program is compiled with the \texttt{--info-heap} flag

The program we will run is an interactive console application ported from \citeauthor*{gebser_asynchronous_2023}~\cite{gebser_asynchronous_2023}. Our code can be found in Appendix \ref{appendix:valgrind_program} and the full output can be found in Appendix \ref{appendix:valgrind_output}.
The program maintains a counter that is incremented every time a clock ticks and can be printed to the screen when a user writes `\textit{show}'. The counter can be manipulated by writing either `\textit{negate}', to negate the current value, or a number $n$ to increase the current value by $n$. 
To implement the desired behaviour, the program uses $clock$ and $console$-input signals, and manipulates them with the $\mathsf{head}$, $\mathsf{tail}$, $map$, $filter$, $switch$, $trigger$, and $interleave$ signal combinators. Through these combinators, the example makes use of Rizzo's core features: signals, $\mathsf{sync}$ construction, $\mathsf{watch}$, $\mathsf{wait}$, $\mathsf{tail}$, and later applications ($\mathsf{laterapp}$). The only Rizzo constructors not used by the example are $\mathsf{never}$ and $\mathsf{ostar}$.

To mimic the program running for an extended period, we adjust the timer so the counter is incremented every 1 millisecond. We then have the program run for 200 seconds with input occasionally being given to the console, which results in about 200,000 steps taken. We run the program with Valgrind's \texttt{--leak-check=full} tool to check for memory leaks and other memory errors. For reproducibility, we made it into a script under \texttt{scripts/valgrind\_evaluation.sh}.

\begin{verbatim}
Steps taken: 198480, Average step time: 154 us, Signals left in heap: 21
==3108== 
==3108== HEAP SUMMARY:
==3108==     in use at exit: 5,811 bytes in 154 blocks
==3108==   total heap usage: 4,169,475 allocs, 4,169,321 frees, 174,719,546 bytes allocated
==3108== 
==3108== LEAK SUMMARY:
==3108==    definitely lost: 0 bytes in 0 blocks
==3108==    indirectly lost: 0 bytes in 0 blocks
==3108==      possibly lost: 0 bytes in 0 blocks
==3108==    still reachable: 5,811 bytes in 154 blocks
==3108==         suppressed: 0 bytes in 0 blocks
\end{verbatim}

The output from Valgrind's leak summary has four categories: definitely lost, indirectly lost, possibly lost, and still reachable. Definitely lost means memory has been allocated but no remaining pointer refers to it, so it is a confirmed memory leak. Indirectly lost means memory is only reachable from another leaked block, so it is leaked as a consequence of the original loss. Possibly lost means Valgrind cannot determine with certainty whether the memory is still accessible, so it may indicate a leak. Still reachable means the memory was not freed before program termination, but valid pointers to it still exist, so it is not considered a true leak. Since the leak summary reports 0 bytes for definitely lost, indirectly lost, and possibly lost, this output shows that the program does not have a memory leak according to Valgrind.

